\documentclass[12pt]{article}

\usepackage{preamble}

% ---------------------------------------------------------------------------
% Document-specific macros
% ---------------------------------------------------------------------------
\newcommand{\anote}[1]{\textcolor{red}{\small\textbf{(Ashwin:} #1\textbf{)}}}
\newcommand{\todo}[1]{\textcolor{\mysubcolor}{\small\textbf{[TODO:} #1\textbf{]}}}

\newcommand{\alg}{\mathcal{A}}             % finite-dimensional C*-algebra
\newcommand{\hilb}{\mathcal{H}}            % Hilbert space
\newcommand{\dil}{\mathcal{K}}             % dilation space
\newcommand{\bdd}[1]{\mathcal{B}\br{#1}}   % bounded operators
\newcommand{\Mn}[1]{\mathrm{M}_{#1}}       % matrix algebra M_n(C)
\newcommand{\algtensor}{\odot}             % algebraic tensor product
\newcommand{\op}{^{\mathrm{op}}}           % opposite algebra
\newcommand{\ip}[2]{\left\langle #1 , #2 \right\rangle}
\newcommand{\ipphi}[2]{\left\langle #1 , #2 \right\rangle_{\varphi}}
\newcommand{\rankone}[2]{\theta_{#1, #2}}  % w |-> <w, #2> #1
\newcommand{\Ltwo}{L^2\br{\alg}}
\newcommand{\vect}[1]{\widehat{#1}}
\newcommand{\Ltn}{\hilb \tensor \Ltwo}
\newcommand{\Fphi}{F_{\varphi}}
\newcommand{\Choi}{C_{\varphi}}
\newcommand{\comm}{\mathcal{C}}            % commutant rho(A)'
\newcommand{\Esc}{E_{\complex 1}}          % expectation onto the scalars
\newcommand{\ran}{\operatorname{ran}}

% ---------------------------------------------------------------------------
\title{\textbf{Stinespring's Theorem via the Choi Matrix:\\
an \texorpdfstring{\(L^2\)}{L2} Route in Finite Dimensions}%
\thanks{Acknowledgements common to all authors.}}

\author{
  Author One\thanks{Affiliation details, address, email.} \\
  Affiliation
  \and
  Author Two\thanks{Affiliation details, address, email.} \\
  U.\ Waterloo
}

\date{\today}

\begin{document}

\maketitle

\begin{abstract}
Let \(\alg\) be a finite-dimensional C\(^*\)-algebra with faithful tracial state
\(\tau\). For every linear map \(\varphi : \alg \to \bdd{\hilb}\), Stinespring's
sesquilinear form on \(\hilb \tensor L^2(\alg,\tau)\) is implemented by a single
operator \(\Fphi\). It lies in the commutant
\(\bdd{\hilb} \tensor \alg\op\) of the standard representation and, in matrix
units, equals the Choi matrix. The correspondence \(\varphi \mapsto \Fphi\) is a
linear bijection onto that commutant, and \(\varphi\) is completely positive
exactly when \(\Fphi \ge 0\). This gives a short, self-contained proof of
Stinespring's theorem~\cite{Stinespring1955} in which the quotient-and-completion
step becomes \(\Fphi^{1/2}\), recovers Choi's theorem, and yields Jordan,
Krein-space and two-sided dilations of non-completely-positive maps. For
completely positive maps, \(\Fphi\) is the Radon--Nikodym derivative of
\(\varphi\) with respect to \(a \mapsto \tau(a)\id\)~\cite{Arveson1969,
BelavkinStaszewski1986,Raginsky2003}. \todo{one sentence on the generalization
beyond finite dimensions, if any.}
\end{abstract}

\tableofcontents

% ===========================================================================
\section{Introduction}
\label{sec:intro}

Stinespring's theorem~\cite{Stinespring1955} writes every completely positive map
\(\varphi\) on a unital C\(^*\)-algebra as \(\varphi(a) = V^\adjoint\pi(a)V\). Its
proof equips \(\alg \algtensor \hilb\) with a \(\varphi\)-dependent form and takes
the separated completion. In finite dimensions, Choi's theorem~\cite{Choi1975}
characterizes complete positivity by positivity of the Choi matrix.

We connect the two. For finite-dimensional \(\alg\), we fix the space
\(\Ltn\) in advance, using the GNS space of a faithful trace as in the theory of
finite von Neumann algebras~\cite{AnantharamanPopa}. All dependence on \(\varphi\)
is carried by one operator \(\Fphi\), which turns out to be the Choi matrix.

\paragraph{Prior work.}
For completely positive \(\varphi\), \(\Fphi\) is the Arveson /
Belavkin--Staszewski Radon--Nikodym derivative~\cite{Arveson1969,
BelavkinStaszewski1986} of \(\varphi\) with respect to the expectation onto the scalars
\(\Esc(a) = \tau(a)\id_{\hilb}\). For full matrix algebras, Raginsky computes this
derivative as a multiple of the Choi--Jamio\l{}kowski
operator~\cite[Theorem~5.1]{Raginsky2003}. What we add is the following:
\begin{itemize}
  \item The argument runs in the other direction. The dilation of the reference
    map is explicit, so the construction proves Stinespring's theorem rather
    than assuming it.
  \item It covers all finite-dimensional C\(^*\)-algebras, with \(\Fphi\) placed in
    \(\bdd{\hilb} \tensor \alg\op\) by the commutation theorem
    \(\lambda(\alg)' = R(\alg\op)\).
  \item It extends to arbitrary linear maps.
\end{itemize}
\todo{Read Belavkin--Staszewski in full to confirm that direct sums and
non-positive maps are not already covered there.}

\paragraph{Results.}
\Cref{thm:linear} constructs \(\Fphi\) for every linear map \(\varphi\) and
identifies \(\varphi \mapsto \Fphi\) as a bijection onto the commutant.
\Cref{thm:cp} gives the dilation of completely positive maps.
\Cref{thm:choi} identifies \(\Fphi\) with the Choi matrix.
\Cref{sec:non-cp} treats non-completely-positive maps.
\Cref{sec:extremal} follows Naimark's \S19~\cite{Naimark1972} to derive the
Radon--Nikodym theorem, pure maps and extreme UCP and CPTP maps from \(\Fphi\). It
recovers the Choi and Kraus characterizations, including versions for direct
sums.
\Cref{sec:further} collects open directions. \todo{Generalization beyond the
finite-dimensional case, e.g.\ to finite von Neumann algebras
(\cref{sec:further}).}

\paragraph{Acknowledgements.}
Personal acknowledgements, grants, etc.

% ===========================================================================
\section{Preliminaries}
\label{sec:prelim}

\paragraph{Conventions.}
Inner products are linear in the \emph{first} argument, and \(\rankone{u}{v}\) is
the operator \(w \mapsto \ip{w}{v}u\). For a linear map \(\varphi\) we write
\(\varphi_n \eqdef \id_{\Mn{n}} \tensor \varphi\), and
\(\varphi^\sharp(a) \eqdef \varphi(a^\adjoint)^\adjoint\). \(S'\) denotes the
commutant of a set of operators. \(\alg\op\) is \(\alg\) with product
\(a\op b\op \eqdef (ba)\op\). The map \(\varphi\) is \emph{completely positive} if
every \(\varphi_n\) is positive, and \emph{Hermitian} if \(\varphi^\sharp = \varphi\).

\begin{theorem}[Stinespring~{\cite[Theorem~1]{Stinespring1955}}]
\label{thm:stinespring}
Let \(\alg\) be a unital C\(^*\)-algebra and \(\varphi : \alg \to \bdd{\hilb}\)
completely positive. There are a Hilbert space \(\dil\), a \(*\)-representation
\(\pi\) of \(\alg\) on \(\dil\), and a bounded \(V : \hilb \to \dil\) with
\(\varphi(a) = V^\adjoint\pi(a)V\) and \(\norm{V}^2 = \norm{\varphi(1)}\).
\end{theorem}
\todo{Check the theorem number and the exact statement against the original
paper.}

Such a triple \((\pi, V, \dil)\) is a \emph{Stinespring dilation}. It is
\emph{minimal} if \(\pi(\alg)V\hilb\) spans a dense subspace. Minimal dilations
are unique up to a unitary intertwining \(\pi\) and \(V\)~\cite{Arveson1969},
\cite[Chapter~4]{Paulsen2002}; this is not in Stinespring's paper. His proof
puts the form
\begin{equation}
\label{eq:form}
  \ipphi{x \tensor a}{y \tensor b} \eqdef \ip{\varphi(b^\adjoint a)x}{y}
\end{equation}
on \(\hilb \algtensor \alg\), shows it is positive semidefinite, and takes the
separated completion. In that completion \(\pi\) is left multiplication and
\(Vx = [x \tensor 1]\).

% ===========================================================================
\section{The Standard Form of a Finite-Dimensional Algebra}
\label{sec:standard}

From now on \(\alg\) is a finite-dimensional C\(^*\)-algebra. By the structure
theorem~\cite[Theorem~III.1.1]{Davidson1996}, \cite{Takesaki1979},
\begin{equation}
\label{eq:structure}
  \alg \cong \bigoplus_{k=1}^r \Mn{n_k}(\complex), \qquad a = (a_1, \dots, a_r).
\end{equation}
\todo{Check the theorem number in Davidson and add a precise reference in
Takesaki, Vol.~I.}
Let \(\iota_k\) and \(p_k\) be the inclusion of, and projection onto, the
\(k\)-th summand, and set \(\varphi_k \eqdef \varphi \compose \iota_k\). Then
\(\varphi = \sum_k \varphi_k \compose p_k\). Since \(\iota_k\) and \(p_k\) are
\(*\)-homomorphisms, \(\varphi\) is completely positive if and only if every
\(\varphi_k\) is.

\paragraph{Traces.}
Every faithful tracial state on \(\alg\) has the form
\(\tau(a) = \sum_k t_k \trace(a_k)\) with \(t_k > 0\) and \(\sum_k t_k n_k = 1\).
This is because a trace on \(\Mn{n}\) is a multiple of \(\trace\). The trace is
unique, equal to \(\frac1n\trace\), exactly when \(\alg\) is a factor. Fix such a
\(\tau\), and write \(c_k \eqdef t_k n_k = \tau(\iota_k(1))\).

\paragraph{The space \(\Ltwo\).}
\(\Ltwo = L^2(\alg, \tau)\) is \(\alg\) with inner product
\(\ip{\vect a}{\vect b} \eqdef \tau(b^\adjoint a)\), and \(\Omega \eqdef \vect 1\).
Since \(\tau\) is faithful and \(\dim \alg < \infty\), no quotient or completion is
needed: \(a \mapsto \vect a\) is a linear bijection. The norm
\(\norm{\vect a}_2 = \tau(a^\adjoint a)^{1/2} \le \norm a\) is equivalent to, but
in general different from, the C\(^*\)-norm; for \(E_{11} \in \Mn n\) it equals
\(n^{-1/2}\). The \emph{standard representation} \(\lambda\) and right
multiplication \(R\) are
\[
  \lambda(a)\vect c \eqdef \vect{ac}, \qquad R_b\vect c \eqdef \vect{cb}.
\]

\begin{lemma}
\label{lem:standard}
\begin{enumerate}[label=(\roman*)]
  \item \(\lambda\) is a faithful unital \(*\)-representation of \(\alg\), and
    \(\lambda(\alg)\Omega = \Ltwo\).
  \item \(b\op \mapsto R_b\) is a faithful unital \(*\)-representation of
    \(\alg\op\).
  \item \(\lambda(\alg)' = R(\alg\op)\).
  \item For any Hilbert space \(\hilb\),
    \((\id_{\hilb} \tensor \lambda(\alg))' = \bdd{\hilb} \tensor R(\alg\op)\), an
    algebraic tensor product that is already a von Neumann algebra.
\end{enumerate}
\end{lemma}
\begin{proof}
(i) We have \(\tau(b^\adjoint ca) = \tau((c^\adjoint b)^\adjoint a)\), so
\(\lambda(c)^\adjoint = \lambda(c^\adjoint)\), and \(\lambda(a)\Omega = \vect a\).
(ii) \(R_bR_c = R_{cb}\), which is multiplicativity on \(\alg\op\). Adjoints use
the trace property:
\(\tau(d^\adjoint cb) = \tau(bd^\adjoint c) = \tau((db^\adjoint)^\adjoint c)\).
(iii) Associativity gives \(\supseteq\). Conversely, if \(T\) commutes with
\(\lambda(\alg)\) and \(T\Omega = \vect d\), then
\(T\vect a = T\lambda(a)\Omega = \lambda(a)\vect d = R_d\vect a\).
(iv) Since \(\dim\Ltwo < \infty\), write \(T = \sum_k S_k \tensor T_k\) with
\(S_k \in \bdd\hilb\) linearly independent. If \(T\) commutes with
\(\id \tensor \lambda(c)\), slicing with functionals dual to the \(S_k\) gives
\([T_k, \lambda(c)] = 0\), so each \(T_k \in R(\alg\op)\) by~(iii).
\end{proof}

\begin{remark}
\label{rem:op}
\begin{enumerate}[label=(\alph*)]
  \item In general \(\bdd{\Ltwo} \not\cong \alg \tensor \alg\op\). The algebra
    generated by \(\lambda(\alg)\) and \(R(\alg\op)\) has commutant the centre of
    \(\alg\). For example, \(\complex \oplus \complex\) gives \(\complex^4\)
    versus \(\Mn 2\). Equality holds exactly for factors, and is not needed.
  \item The blockwise transpose \(b\op \mapsto (b_k^\transpose)_k\) is a
    \(*\)-isomorphism \(\alg\op \cong \alg\). So
    \(\bdd\hilb \tensor R(\alg\op) \cong \alg \tensor \bdd\hilb\), non-canonically
    since it depends on the matrix units.
  \item (Blocks.) The map
    \(U : \vect a \mapsto (c_k^{1/2}\vect{a_k})_k\) is a unitary
    \(\Ltwo \to \bigoplus_k L^2(\Mn{n_k}, \frac{1}{n_k}\trace)\). It carries
    \(\lambda\), \(R\) and \(\Omega\) to the direct sums of their counterparts, with
    \(U\Omega = (c_k^{1/2}\Omega_k)_k\). This follows from
    \(\tau(b^\adjoint a) = \sum_k c_k\,\frac{1}{n_k}\trace(b_k^\adjoint a_k)\).
\end{enumerate}
\end{remark}

% ===========================================================================
\section{The Operator \texorpdfstring{\(\Fphi\)}{F}}
\label{sec:F}

Let \(\hilb\) be a Hilbert space, \(\varphi : \alg \to \bdd\hilb\) linear,
\[
  \rho(a) \eqdef \id_{\hilb} \tensor \lambda(a), \qquad
  Vx \eqdef x \tensor \Omega, \qquad
  \comm \eqdef \rho(\alg)' = \bdd\hilb \tensor R(\alg\op).
\]
Since \(\dim\Ltwo < \infty\), \(\Ltn = \hilb \algtensor \alg\), so the
form~\eqref{eq:form} is defined on all of \(\Ltn\). \(V\) is an isometry, but
\(V^\adjoint\rho(a)V = \tau(a)\id_{\hilb}\). Thus \((\rho, V, \Ltn)\) is the
minimal dilation of the expectation onto the scalars \(\Esc(a) = \tau(a)\id\), not of \(\varphi\);
the map \(\varphi\) enters only through~\eqref{eq:form}.

\begin{theorem}[Linear maps]
\label{thm:linear}
\begin{enumerate}[label=(\alph*)]
  \item \label{it:exists}
    There is a unique \(\Fphi \in \bdd{\Ltn}\) with
    \(\ipphi uv = \ip{\Fphi u}{v}\) for all \(u, v\).
  \item \label{it:comm}
    \(\Fphi \in \comm\), and
    \(\varphi(a) = V^\adjoint\rho(a)\Fphi V\) for all \(a \in \alg\).
  \item \label{it:bij}
    \(\varphi \mapsto \Fphi\) is a linear bijection from the linear maps
    \(\alg \to \bdd\hilb\) onto \(\comm\), with inverse
    \(T \mapsto \varphi_T(a) \eqdef V^\adjoint\rho(a)TV\).
  \item \label{it:adj}
    \(F_{\varphi^\sharp} = \Fphi^\adjoint\). In particular, \(\varphi\) is
    Hermitian if and only if \(\Fphi = \Fphi^\adjoint\).
\end{enumerate}
\end{theorem}
\begin{proof}
Fix an orthonormal basis \(\set{\vect{u}_p}_{p=1}^m\) of \(\Ltwo\). Write
\(u = \sum_p x_p \tensor \vect u_p\) and \(v = \sum_q y_q \tensor \vect u_q\), so
that \(\norm u^2 = \sum\norm{x_p}^2\).

(a) The form is induced by a map linear in \(x, a\) and conjugate linear in
\(y, b\), so it is well defined. It is bounded:
\(\abs{\ipphi uv} \le m\max_{p,q}\norm{\varphi(u_q^\adjoint u_p)}\,\norm u\norm v\).
A bounded sesquilinear form is implemented by a unique bounded
operator~\cite[Theorem~II.2.2]{Conway1990}. \todo{Check the theorem number in
Conway.}

(b) Using \(\rho(c)^\adjoint = \rho(c^\adjoint)\),
\[
  \ip{\Fphi\rho(c)(x \tensor \vect a)}{y \tensor \vect b}
  = \ip{\varphi(b^\adjoint ca)x}{y}
  = \ip{\Fphi(x \tensor \vect a)}{\rho(c^\adjoint)(y \tensor \vect b)},
\]
so \(\Fphi\rho(c) = \rho(c)\Fphi\), and \(\comm = \bdd\hilb \tensor R(\alg\op)\)
by \cref{lem:standard}(iv). Then
\(\ip{V^\adjoint\rho(a)\Fphi Vx}{y} = \ip{\Fphi(x \tensor \vect a)}{y \tensor \Omega}
= \ip{\varphi(a)x}{y}\).

(c) Linearity follows from uniqueness in~(a), and injectivity from~(b). For
\(T \in \comm\),
\(\ip{\varphi_T(b^\adjoint a)x}{y}
= \ip{\rho(b^\adjoint a)T(x \tensor \Omega)}{y \tensor \Omega}
= \ip{T(x \tensor \vect a)}{y \tensor \vect b}\), so \(F_{\varphi_T} = T\).

(d) \(\ip{\Fphi^\adjoint(x \tensor \vect a)}{y \tensor \vect b}
= \overline{\ip{\varphi(a^\adjoint b)y}{x}}
= \ip{\varphi^\sharp(b^\adjoint a)x}{y}\). Then apply~(c).
\end{proof}

\begin{remark}
\label{rem:rn}
For completely positive \(\varphi\), \cref{thm:linear}\ref{it:comm} is the
Radon--Nikodym theorem~\cite{Arveson1969,BelavkinStaszewski1986} for \(\varphi\)
relative to \(\Esc\), whose minimal dilation is \((\rho, V, \Ltn)\).
\Cref{thm:linear} needs no positivity, and its proof does not use
\cref{thm:stinespring}. For \(\hilb = \complex^d\), both sides of
\ref{it:bij} have dimension \(d^2\dim\alg\).
\end{remark}

% ===========================================================================
\section{Completely Positive Maps}
\label{sec:cp}

\begin{theorem}
\label{thm:cp}
\begin{enumerate}[label=(\alph*)]
  \item \label{it:pos}
    \(\varphi\) is completely positive if and only if \(\Fphi \ge 0\).
  \item \label{it:dil}
    If \(\Fphi \ge 0\), let \(\dil \eqdef \overline{\ran}\,\Fphi^{1/2}\),
    \(\pi \eqdef \rho|_{\dil}\) and \(W \eqdef \Fphi^{1/2}V\). Then \((\pi, W, \dil)\)
    is a minimal Stinespring dilation of \(\varphi\) with \(W^\adjoint W = \varphi(1)\).
  \item \label{it:compare}
    The map \([x \tensor a] \mapsto \Fphi^{1/2}(x \tensor \vect a)\) is a unitary
    from Stinespring's space onto \(\dil\), intertwining the representations and
    the dilation operators.
\end{enumerate}
\end{theorem}
\begin{proof}
(b) \(\Fphi^{1/2} \in \comm\), so \(\dil\) reduces \(\rho\) and \(\pi\) is a unital
\(*\)-representation. By \cref{thm:linear}\ref{it:comm},
\(W^\adjoint\pi(a)W = V^\adjoint\rho(a)\Fphi V = \varphi(a)\). Minimality holds
because \(\Span\pi(\alg)W\hilb = \Fphi^{1/2}(\Ltn)\).

(a) Suppose \(\varphi\) is completely positive. For
\(u = \sum_p x_p \tensor \vect{a_p}\), \(\ip{\Fphi u}{u} =
\ip{\varphi_m([a_q^\adjoint a_p]_{q,p})\vec x}{\vec x} \ge 0\), since
\([a_q^\adjoint a_p] = X^\adjoint X\) for the row \(X = (a_1, \dots, a_m)\). This is
Stinespring's argument~\cite{Stinespring1955}; see also
\cite[Theorem~4.1]{Paulsen2002}. Conversely, if \(\Fphi \ge 0\) then~(b), which
used only \(\Fphi \ge 0\), exhibits \(\varphi\) as a compression of a
\(*\)-representation, which is completely positive.

(c) \(\norm{\Fphi^{1/2}u}^2 = \ipphi uu\), so the map is isometric on the
separated quotient. Its range is dense, and it intertwines because
\(\Fphi^{1/2} \in \comm\) and \([x \tensor 1] \mapsto Wx\).
\end{proof}

By \cref{rem:op}(c), everything splits over the summands of~\eqref{eq:structure}.
Under \(\id \tensor U\), \(\Fphi \mapsto \bigoplus_k c_k^{-1}F_{\varphi_k}\),
\(\dil = \bigoplus_k \dil_k\), \(\pi = \bigoplus_k \pi_k\) and
\(W = (W_k)_k\), where \(F_{\varphi_k}\) is taken with respect to
\(\frac{1}{n_k}\trace\). This follows from
\(\ipphi{x \tensor \vect a}{y \tensor \vect b} = \sum_k
\ip{x \tensor \vect{a_k}}{y \tensor \vect{b_k}}_{\varphi_k}\) and uniqueness in
\cref{thm:linear}\ref{it:exists}. So it suffices to treat a single factor
\(\Mn n\).

\subsection{Minimality and cyclicity: the GNS case}
\label{sec:gns}

Stinespring's theorem generalizes the GNS construction, and minimality of a
dilation generalizes cyclicity of the GNS vector. Let \((\pi, W, \dil)\) be any
Stinespring dilation of a completely positive \(\varphi : \alg \to \bdd\hilb\).
For a unit vector \(x \in \hilb\), the vector \(Wx\) satisfies
\(\ip{\pi(a)Wx}{Wx} = \ip{\varphi(a)x}{x}\). So the cyclic subspace
\(\overline{\pi(\alg)Wx}\) with the vector \(Wx\) is a GNS representation of
the positive functional \(\omega_x \compose \varphi\), where
\(\omega_x = \ip{\cdot\,x}{x}\). The dilation is minimal exactly when \(\dil\) is
the closed linear span of these GNS spaces as \(x\) ranges over \(\hilb\).
Two cases are worth recording.
\begin{enumerate}[label=(\roman*)]
  \item \(\hilb = \complex\). Then \(W\) is determined by \(\xi \eqdef W1\), and
    \(\varphi(a) = \ip{\pi(a)\xi}{\xi}\). Minimality says \(\pi(\alg)\xi\) is
    dense, which is cyclicity of \(\xi\). So a minimal Stinespring dilation of a
    positive functional is exactly a GNS triple. Uniqueness of minimal dilations
    up to unitary equivalence~\cite[Proposition~4.2]{Paulsen2002} specializes to
    uniqueness of the GNS representation. \todo{Check the proposition number in
    Paulsen.}
  \item The model of \cref{sec:F}. By \cref{lem:standard}(i), the vector
    \(\Omega\) is cyclic for \(\lambda\), so
    \(\Span\rho(\alg)V\hilb = \Ltn\). Thus \(\Ltwo\) is the GNS space of \(\tau\)
    and \((\rho, V, \Ltn)\) is the minimal dilation of \(\Esc\). The minimality in
    \cref{thm:cp}\ref{it:dil} is inherited from this one GNS cyclicity, because
    \(\Fphi^{1/2}\) commutes with \(\rho(\alg)\):
    \[
      \Span\pi(\alg)W\hilb = \Fphi^{1/2}\Span\rho(\alg)V\hilb .
    \]
\end{enumerate}

For \(\hilb = \complex\) the construction of \cref{sec:F} becomes the familiar
description of the GNS space of a functional inside \(L^2(\alg,\tau)\). Here
\(\Ltn = \Ltwo\), \(\rho = \lambda\), \(V1 = \Omega\) and \(\comm = R(\alg\op)\).
Since \(\tau\) is faithful and \(\dim\alg < \infty\), every linear functional
has the form \(f = \tau(\cdot\,h_f)\) for a unique \(h_f \in \alg\), the
\emph{density} of \(f\) relative to \(\tau\).

\begin{proposition}[Functionals]
\label{prop:gns}
Let \(f = \tau(\cdot\,h)\) be a linear functional on \(\alg\).
\begin{enumerate}[label=(\alph*)]
  \item \label{it:gns-F}
    \(F_f = R_h\).
  \item \label{it:gns-pos}
    \(f\) is positive \(\iff\) \(h \ge 0\) \(\iff\) \(F_f \ge 0\) \(\iff\) \(f\) is
    completely positive.
  \item \label{it:gns-dil}
    If \(f \ge 0\), let \(p\) be the support projection of \(h\). The dilation of
    \cref{thm:cp}\ref{it:dil} is
    \[
      \dil = \set{\vect{ap} : a \in \alg}, \qquad
      \pi(a) = \lambda(a)|_{\dil}, \qquad
      \xi_f \eqdef W1 = \vect{h^{1/2}},
    \]
    and \(\vect a \mapsto \vect{ah^{1/2}}\) induces a unitary from the GNS space
    \(L^2(\alg, f)\) onto \(\dil\) that carries the GNS vector \(\vect 1\) to
    \(\xi_f\). The GNS null space is
    \(\set{a : f(a^\adjoint a) = 0} = \alg(1 - p)\).
  \item \label{it:gns-faithful}
    \(f\) is faithful \(\iff\) \(h\) is invertible \(\iff\) \(\dil = \Ltwo\). In
    that case \(\xi_f\) is cyclic and separating for \(\lambda(\alg)\).
  \item \label{it:gns-choi}
    For \(\alg = \Mn n\) with \(\tau = \frac1n\trace\), the density matrix of
    \(f\) is \(D_f = h/n\), and its Choi matrix is
    \(C_f = [f(E_{ij})]_{ij} = D_f^\transpose\). Under the identification of
    \cref{thm:choi}, \(F_f = R_h\) corresponds to
    \(n(\id_n \tensor C_f) = \id_n \tensor h^\transpose\).
\end{enumerate}
\end{proposition}
\begin{proof}
(a) \(\ip{R_h\vect a}{\vect b} = \tau(b^\adjoint ah) = f(b^\adjoint a)
= \ip{F_f\vect a}{\vect b}\).
(b) If \(f\) is positive and \(e\) is a unit vector in the \(k\)-th summand, then
for the rank-one projection \(q = \iota_k(\rankone ee)\) we get
\(0 \le f(q) = \tau(qhq) = t_k\ip{h_ke}{e}\). So \(h \ge 0\). Since
\(b\op \mapsto R_b\) is a \(*\)-representation of \(\alg\op\) and \(h\op \ge 0\) in
\(\alg\op\), \(h \ge 0\) gives \(R_h \ge 0\). Then \(F_f \ge 0\) by~(a), and
\(f\) is completely positive by \cref{thm:cp}\ref{it:pos}. Completely positive
maps are positive.
(c) \(F_f^{1/2} = R_{h^{1/2}}\), whose range is
\(\set{\vect{ah^{1/2}}} = \set{\vect{ap}}\), since \(h^{1/2}\) is invertible in
\(p\alg p\). Also \(W1 = R_{h^{1/2}}\Omega = \vect{h^{1/2}}\). Finally
\(\norm{\vect{ah^{1/2}}}_2^2 = \tau(h^{1/2}a^\adjoint ah^{1/2}) = f(a^\adjoint a)\),
so the map is isometric on the GNS quotient. It kills exactly
\(\set{a : ah^{1/2} = 0} = \alg(1-p)\), it is onto \(\dil\), and it intertwines
left multiplication.
(d) Faithfulness means the null space \(\alg(1-p)\) is zero, that is \(p = 1\).
Then \(\pi(a)\xi_f = 0\) forces \(ah^{1/2} = 0\), so \(a = 0\).
(e) \(f(E_{ij}) = \frac1n\trace(E_{ij}h) = \frac1n h_{ji}\).
\end{proof}

\begin{remark}
In \ref{it:gns-dil}, \(\xi_f = \vect{h^{1/2}}\) is the representative of \(f\) in
the positive cone \(\set{\vect b : b \ge 0}\) of \(\Ltwo\). For faithful \(f\)
this is the standard-form picture of Tomita--Takesaki theory, here with no
modular theory needed because \(\tau\) is a trace. \todo{Add a reference for the
natural cone and vector representatives (e.g.\ Takesaki, Vol.~II).}
\end{remark}

\paragraph{Abelian codomains.}
Positive functionals are automatically completely positive, as
\ref{it:gns-pos} shows. More generally, a positive map into an abelian
C\(^*\)-algebra is completely positive~\cite[Theorem~4]{Stinespring1955},
\cite[Theorem~3.9]{Paulsen2002}. \todo{Check both theorem numbers.} In finite
dimensions this also follows from the operator \(\Fphi\), and the proof gives
the minimal dilation explicitly as a direct sum of amplified GNS
representations.

\begin{corollary}
\label{cor:abelian}
Let \(\hilb\) be finite-dimensional, and let \(\varphi : \alg \to \bdd\hilb\) be
positive with range in an abelian \(*\)-subalgebra. Let
\(P_1, \dots, P_m\) be the minimal projections of the unital \(*\)-algebra this
range generates. Then
\(\varphi(a) = \sum_k f_k(a)P_k\) for positive functionals
\(f_k = \tau(\cdot\,h_k)\), and
\[
  \Fphi = \sum_{k=1}^m P_k \tensor R_{h_k} \ge 0 .
\]
So \(\varphi\) is completely positive. Its minimal dilation is
\(\dil = \bigoplus_k P_k\hilb \tensor \dil_{f_k}\) with
\(\pi = \bigoplus_k \id \tensor \pi_{f_k}\) and
\(Wx = \sum_k P_kx \tensor \xi_{f_k}\), where \((\pi_{f_k}, \xi_{f_k}, \dil_{f_k})\)
is the GNS triple of \cref{prop:gns}\ref{it:gns-dil}.
\end{corollary}
\begin{proof}
The algebra generated by the range is spanned by the pairwise orthogonal
projections \(P_k\), which sum to \(\id\), so \(\varphi(a) = \sum_k f_k(a)P_k\)
for linear functionals \(f_k\). For a unit vector \(x_k \in P_k\hilb\),
\(f_k(a) = \ip{\varphi(a)x_k}{x_k}\), so each \(f_k\) is positive. Then
\[
  \ip{\varphi(b^\adjoint a)x}{y}
  = \sum_k f_k(b^\adjoint a)\ip{P_kx}{y}
  = \sum_k \ip{(P_k \tensor R_{h_k})(x \tensor \vect a)}{y \tensor \vect b},
\]
by \cref{prop:gns}\ref{it:gns-F}. Each \(P_k \tensor R_{h_k}\) is a product of
commuting positive operators, so \(\Fphi \ge 0\), and \cref{thm:cp}\ref{it:pos}
applies. Finally,
\(\Fphi^{1/2} = \sum_k P_k \tensor R_{h_k^{1/2}}\), and
\cref{prop:gns}\ref{it:gns-dil} identifies its range and \(\Fphi^{1/2}V\).
\end{proof}
\todo{For infinite-dimensional \(\hilb\) and a general abelian von Neumann
algebra, replace the sum by a spectral integral
\(\Fphi = \int dE(\omega) \tensor R_{h(\omega)}\), and compare with Stinespring's
proof. Note also the dual statement: a positive map \emph{on} an abelian
algebra is completely positive~\cite{Stinespring1955}, which here would say
that for abelian \(\alg\), \(\Fphi \ge 0\) follows from positivity alone.}

% ===========================================================================
\section{The Choi Matrix}
\label{sec:choi}

\begin{theorem}
\label{thm:choi}
Let \(\alg = \Mn n\), \(\tau = \frac1n\trace\), \(\varphi : \Mn n \to \bdd\hilb\)
linear, and \(\Choi \eqdef \sum_{j,l} E_{jl} \tensor \varphi(E_{jl})\). Under
\(\vect{E_{ij}} \mapsto n^{-1/2}e_i \tensor e_j\), and with \(\hilb\) moved last,
\[
  \Fphi = n\,(\id_n \tensor \Choi).
\]
For \(\alg = \bigoplus_k \Mn{n_k}\), accordingly
\(\Fphi \cong \bigoplus_k t_k^{-1}(\id_{n_k} \tensor C_{\varphi_k})\).
\end{theorem}
\begin{proof}
In any orthonormal basis, comparing matrix entries gives
\(\Fphi = \sum_{p,q}\varphi(u_p^\adjoint u_q) \tensor \rankone{\vect u_p}{\vect u_q}\).
Take \(u_{ij} = n^{1/2}E_{ij}\). Then
\(u_{ij}^\adjoint u_{i'l} = n\delta_{ii'}E_{jl}\), and
\(\sum_i \rankone{\vect u_{ij}}{\vect u_{il}} \mapsto \id_n \tensor E_{jl}\).
The direct-sum case follows from \cref{sec:cp} and \(c_k^{-1}n_k = t_k^{-1}\).
\end{proof}

\begin{corollary}[Choi~\cite{Choi1975}]
\label{cor:choi}
\(\varphi : \Mn n \to \bdd\hilb\) is completely positive if and only if
\(\Choi \ge 0\). For \(\alg = \bigoplus_k \Mn{n_k}\), \(\varphi\) is completely
positive if and only if every \(C_{\varphi_k} \ge 0\).
\end{corollary}

\begin{remark}
\label{rem:legs}
Under this identification, \(\lambda(a) = a \tensor \id_n\) and
\(R_b = \id_n \tensor b^\transpose\). So \(\id_n\) is the \(\lambda\)-leg and the
Choi matrix sits on the \(\alg\op\)-leg, as \cref{thm:linear}\ref{it:comm}
predicts. In \(\bdd\hilb \tensor \alg\op\),
\(\Fphi = n\sum_{j,l}\varphi(E_{jl}) \tensor E_{lj}\op\).
For \(\hilb = \complex^d\), \(\dim\dil = n\operatorname{rank}\Choi\). This matches
\(\dil \cong \complex^n \tensor \complex^\kappa\), \(\pi(a) = a \tensor \id_\kappa\),
with multiplicity \(\kappa = \operatorname{rank}\Choi\). The Choi identification
holds for every linear \(\varphi\), since its proof uses no positivity.
\end{remark}

% ===========================================================================
\section{Maps That Are Not Completely Positive}
\label{sec:non-cp}

\begin{theorem}
\label{thm:non-cp}
Let \(\varphi : \alg \to \bdd\hilb\) be linear.
\begin{enumerate}[label=(\alph*)]
  \item \label{it:jordan}
    (Jordan.) If \(\varphi\) is Hermitian and \(\Fphi = F_+ - F_-\) with
    \(F_\pm \ge 0\) and \(F_+F_- = 0\), then \(F_\pm \in \comm\). Moreover
    \(\varphi = \varphi_{F_+} - \varphi_{F_-}\), a difference of completely
    positive maps.
  \item \label{it:krein}
    (Krein.) If \(\varphi\) is Hermitian, let \(\dil = \overline{\ran}\abs\Fphi\),
    \(\pi = \rho|_{\dil}\), \(W = \abs\Fphi^{1/2}V\) and
    \(\Sigma = \operatorname{sgn}(\Fphi)|_{\dil}\). Then
    \(\Sigma = \Sigma^\adjoint = \Sigma^{-1}\) commutes with \(\pi(\alg)\), and
    \(\varphi(a) = W^\adjoint\pi(a)\Sigma W\).
  \item \label{it:four}
    (Four maps.) \(\varphi = \varphi_1 + i\varphi_2\) with
    \(\varphi_1 = \frac12(\varphi + \varphi^\sharp)\) and
    \(\varphi_2 = \frac1{2i}(\varphi - \varphi^\sharp)\) Hermitian, and
    \(F_{\varphi_1} = \operatorname{Re}\Fphi\),
    \(F_{\varphi_2} = \operatorname{Im}\Fphi\). Hence \(\varphi\) is a combination
    of four completely positive maps.
  \item \label{it:two-sided}
    (Two-sided.) With polar decomposition \(\Fphi = U\abs\Fphi\),
    \(W_1 = \abs{\Fphi^\adjoint}^{1/2}V\) and \(W_2 = U\abs\Fphi^{1/2}V\), we have
    \(\varphi(a) = W_1^\adjoint\rho(a)W_2\) and \(\norm{W_1}\norm{W_2} \le \norm\Fphi\).
\end{enumerate}
\end{theorem}
\begin{proof}
All of \(F_\pm\), \(\abs\Fphi\), \(\operatorname{sgn}\Fphi\), \(U\) and
\(\abs{\Fphi^\adjoint}\) are obtained from \(\Fphi\) by functional calculus or polar
decomposition, so they lie in the von Neumann algebra \(\comm\). (a) follows from
\cref{thm:linear}\ref{it:bij} and \cref{thm:cp}\ref{it:pos}.
(b) follows from \(\Fphi = \abs\Fphi^{1/2}\operatorname{sgn}(\Fphi)\abs\Fphi^{1/2}\) and
\cref{thm:linear}\ref{it:comm}. (c) follows from \cref{thm:linear}\ref{it:adj}
and~(a). For (d), use
\(\Fphi = \abs{\Fphi^\adjoint}^{1/2}U\abs\Fphi^{1/2}\) and
\cref{thm:linear}\ref{it:comm}.
\end{proof}

\begin{remark}
In~\ref{it:krein}, \(\dil\) with \(\ip{\Sigma\cdot}{\cdot}\) is a Krein space with fundamental symmetry
\(\Sigma\), and \(\pi(a^\adjoint)\) is the \(\Sigma\)-adjoint of \(\pi(a)\). (We
avoid the customary letter \(J\), the standard notation for the conjugation
\(\hat a \mapsto \widehat{a^\adjoint}\) on \(\Ltwo\).) For completely positive \(\varphi\), \(\Sigma = \id\) and this is \cref{thm:cp}\ref{it:dil}.
\todo{Locate the existing literature on Krein-space Stinespring dilations of
Hermitian maps and compare.}

Part~\ref{it:two-sided} has the shape of the Wittstock / Haagerup--Paulsen
theorem~\cite[Theorem~8.4]{Paulsen2002} \todo{check theorem number}, where
\(\norm{V_1}\norm{V_2} = \norm\varphi_{cb}\) can be attained.
\todo{Compare \(\norm\Fphi\) with \(\norm\varphi_{cb}\). By \cref{thm:choi}, for
\(\alg = \Mn n\), \(\norm\Fphi = n\norm\Choi\). Can the polar decomposition be
chosen to attain \(\norm\varphi_{cb}\)?}

For positive maps that are not completely positive, \(\Fphi\) is self-adjoint but
\(F_- \ne 0\). \todo{Characterize positivity, or \(k\)-positivity, of \(\varphi\)
in terms of \(\Fphi\). For \(\alg = \Mn n\) this should recover block-positivity
of the Choi matrix.}
\end{remark}

% ===========================================================================
\section{Domination and Extremal Maps}
\label{sec:extremal}

This section follows the structure of Naimark's treatment of positive functionals
and irreducible representations~\cite[\S19]{Naimark1972}, with positive functionals
replaced by CP maps and the GNS form replaced by \(\Fphi\). The dictionary is:
\begin{center}
\small
\begin{tabular}{p{0.44\textwidth}p{0.48\textwidth}}
  Naimark \S19 & here \\ \hline
  positive functional \(f\), cyclic representation \(x \mapsto A_x\)
    & CP map \(\varphi\), dilation \((\pi, W, \dil)\) of \cref{thm:cp} \\
  \(f_1\) dominated by \(f\): \(\lambda f - f_1 \ge 0\)
    & \(\psi \le \lambda\varphi\): \(\lambda\varphi - \psi\) CP \\
  \(f_1(x) = (A_x B\xi_0, \xi_0)\), \(B \in \{A_x\}'\) (Thm.~1)
    & \(\psi(a) = W^\adjoint\pi(a)TW\), \(T \in \pi(\alg)'\) (\cref{thm:rn}) \\
  \(C_f \cong R_f'\), \(f \leftrightarrow \id\) (subsec.~2)
    & \(\mathcal{D}_\varphi \cong \pi(\alg)'\), \(\varphi \leftrightarrow \id\) (\cref{thm:rn}) \\
  indecomposable \(\Leftrightarrow\) irreducible (Thm.~2)
    & pure \(\Leftrightarrow\) irreducible \(\Leftrightarrow\) Choi rank one (\cref{thm:pure}) \\
  normalized states, Krein--Milman (subsec.~4)
    & UCP / CPTP maps and their extreme points (\cref{thm:extreme}) \\
\end{tabular}
\end{center}
\todo{Check the edition and page numbers of Naimark: the attached scan is
\S19, pp.~266--270.}

Throughout, \(\varphi : \alg \to \bdd\hilb\) is completely positive,
\((\pi, W, \dil)\) is its dilation from \cref{thm:cp}, and \(P\) is the projection
onto \(\dil = \overline{\ran}\,\Fphi^{1/2}\). Since \(P\) is the support projection
of \(\Fphi \in \comm\), we have \(P \in \comm\). So an operator \(T\) on \(\dil\)
commutes with \(\pi(\alg)\) exactly when \(T \oplus 0 \in \comm\), and we identify
\(\pi(\alg)'\) with \(P\comm P\). For \(T \in \pi(\alg)'\) put
\[
  \psi_T(a) \eqdef W^\adjoint \pi(a) T W .
\]
For maps \(\psi, \varphi\) we write \(\psi \le \varphi\) if \(\varphi - \psi\) is
completely positive.

\subsection{Domination (Naimark \S19.1--2)}
\label{sec:domination}

\begin{theorem}[Radon--Nikodym]
\label{thm:rn}
\begin{enumerate}[label=(\alph*)]
  \item \label{it:rn-order}
    For linear \(\psi\) and \(\lambda \ge 0\): \(0 \le \psi \le \lambda\varphi\) if
    and only if \(0 \le F_\psi \le \lambda\Fphi\).
  \item \label{it:rn-T}
    If \(0 \le \psi \le \lambda\varphi\), there is a unique \(T \in \pi(\alg)'\) with
    \(0 \le T \le \lambda\) and \(\psi = \psi_T\). Moreover
    \(F_{\psi_T} = \Fphi^{1/2}\,T\,\Fphi^{1/2}\).
  \item \label{it:rn-alg}
    \(T \mapsto \psi_T\) is a linear bijection from \(\pi(\alg)'\) onto
    \(\mathcal{D}_\varphi \eqdef \Span\set{\psi : 0 \le \psi \le \lambda\varphi
    \text{ for some } \lambda}\). It satisfies \(\psi_{T^\adjoint} = \psi_T^\sharp\),
    \(\psi_{\id} = \varphi\), and \(\psi_T\) is CP exactly when \(T \ge 0\).
\end{enumerate}
\end{theorem}
\begin{proof}
(a) By \cref{thm:linear}\ref{it:bij}, \(F\) is linear, and by
\cref{thm:cp}\ref{it:pos}, \(\lambda\varphi - \psi\) is CP if and only if
\(F_{\lambda\varphi - \psi} = \lambda\Fphi - F_\psi \ge 0\).

(b) We follow Naimark's proof of~\cite[\S19, Theorem~1]{Naimark1972}. On the dense
subspace \(\Fphi^{1/2}(\Ltn)\) of \(\dil\), define
\[
  \br{\Fphi^{1/2}u, \Fphi^{1/2}v}_1 \eqdef \ip{F_\psi u}{v}.
\]
It is well defined. If \(\Fphi^{1/2}u = 0\) then
\(\ip{F_\psi u}{u} \le \lambda\norm{\Fphi^{1/2}u}^2 = 0\), so \(F_\psi^{1/2}u = 0\)
and \(F_\psi u = 0\). The same inequality gives
\(0 \le (\xi, \xi)_1 \le \lambda\norm\xi^2\). So the form extends to \(\dil\) and is
implemented by \(T \in \bdd\dil\) with \(0 \le T \le \lambda\) and
\(F_\psi = \Fphi^{1/2}T\Fphi^{1/2}\). For \(c \in \alg\), using
\(\Fphi^{1/2}, F_\psi \in \comm\),
\[
  \ip{T\pi(c)\Fphi^{1/2}u}{\Fphi^{1/2}v}
  = \ip{F_\psi\rho(c)u}{v}
  = \ip{F_\psi u}{\rho(c^\adjoint)v}
  = \ip{\pi(c)T\Fphi^{1/2}u}{\Fphi^{1/2}v},
\]
so \(T \in \pi(\alg)'\). This is Naimark's equation~(3). Next,
\[
  \psi_T(a) = V^\adjoint\Fphi^{1/2}\rho(a)T\Fphi^{1/2}V
  = V^\adjoint\rho(a)\br{\Fphi^{1/2}T\Fphi^{1/2}}V,
  \qquad \Fphi^{1/2}T\Fphi^{1/2} \in \comm .
\]
So
\(F_{\psi_T} = \Fphi^{1/2}T\Fphi^{1/2} = F_\psi\) by \cref{thm:linear}\ref{it:bij},
hence \(\psi_T = \psi\). For uniqueness: \(\Fphi^{1/2}T\Fphi^{1/2} = 0\) forces
\(T = 0\), because \(\ran\Fphi^{1/2}\) is dense in \(\dil\).

(c) The formula \(F_{\psi_T} = \Fphi^{1/2}T\Fphi^{1/2}\) holds for every
\(T \in \pi(\alg)'\), so \(T \mapsto \psi_T\) is linear and injective. If
\(T \ge 0\) then \(\psi_T(a) = (T^{1/2}W)^\adjoint\pi(a)(T^{1/2}W)\) is CP and
\(\psi_T \le \norm T\varphi\). Every element of the C\(^*\)-algebra
\(\pi(\alg)'\) is a combination of four positive ones, so the image is
\(\mathcal{D}_\varphi\) by~(b). Finally,
\(\psi_T^\sharp(a) = (W^\adjoint\pi(a^\adjoint)TW)^\adjoint = W^\adjoint T^\adjoint\pi(a)W
= \psi_{T^\adjoint}(a)\). If \(\psi_T\) is CP, then
\(F_{\psi_T} = \Fphi^{1/2}T\Fphi^{1/2} \ge 0\), and so \(T \ge 0\) by density of
\(\ran\Fphi^{1/2}\).
\end{proof}

\begin{remark}[The expectation onto the scalars as the universal dominant]
\label{rem:universal}
Let \(\Esc(a) \eqdef \tau(a)\id_{\hilb}\) be the \(\tau\)-preserving expectation
onto the scalars, viewed as a map into \(\bdd{\hilb}\). For \(\alg = \Mn{n}\) and
\(\hilb = \complex^n\) it is the completely depolarizing channel, written
\(\Omega\) in~\cite{Watrous2018}. We have \(F_{\Esc} = \id\), so its dilation is
\((\rho, V, \Ltn)\) and \(\pi(\alg)' = \comm\). Every CP \(\varphi\) satisfies
\(\varphi \le \norm\Fphi\,\Esc\) by \cref{thm:rn}\ref{it:rn-order}, so
\(\mathcal{D}_{\Esc}\) is the space of all linear maps. For \(\Esc\),
\cref{thm:rn}\ref{it:rn-alg} is exactly \cref{thm:linear}\ref{it:bij}. So
\cref{sec:F} is Naimark's \S19.2 for the reference map \(\Esc\), which plays the
role of \(f\) and corresponds to the identity operator. For general
\(\varphi\), \cref{thm:rn} is Arveson's Radon--Nikodym
theorem~\cite[Section~1.4]{Arveson1969}, here proved without assuming a
Stinespring dilation. \todo{Check the theorem numbers in Arveson's Section~1.4.}
\end{remark}

\subsection{Pure maps (Naimark \S19.3)}
\label{sec:pure}

A nonzero CP map \(\varphi\) is \emph{pure} if every CP \(\psi \le \varphi\) is a
scalar multiple of \(\varphi\), that is, if \(\varphi\) spans an extreme ray of the
cone of CP maps. In the matrix case we use the following dictionary between the
Choi matrix and Kraus operators.

\begin{lemma}[Kraus dictionary]
\label{lem:kraus}
Let \(\alg = \Mn n\). For bounded \(V_1, V_2 : \hilb \to \complex^n\), put
\(v_i \eqdef \sum_j e_j \tensor V_i^\adjoint e_j \in \complex^n \tensor \hilb\).
Then the Choi matrix of \(a \mapsto V_1^\adjoint a V_2\) is \(\rankone{v_1}{v_2}\).
Every \(v \in \complex^n \tensor \hilb\) arises this way from a unique \(V\),
and \(V_1, \dots, V_\kappa\) are linearly independent if and only if
\(v_1, \dots, v_\kappa\) are. In particular, if \(\Choi \ge 0\) has finite rank
\(\kappa\), its spectral decomposition \(\Choi = \sum_{i=1}^\kappa \rankone{v_i}{v_i}\)
gives a Kraus decomposition \(\varphi(a) = \sum_{i=1}^\kappa V_i^\adjoint a V_i\) with
\(V_1, \dots, V_\kappa\) linearly independent, and \(\ran\Choi = \Span\set{v_i}\).
\end{lemma}
\begin{proof}
We have \(V_1^\adjoint E_{jl}V_2 = \rankone{V_1^\adjoint e_j}{V_2^\adjoint e_l}\), so
\(\sum_{j,l} E_{jl} \tensor V_1^\adjoint E_{jl}V_2 = \rankone{v_1}{v_2}\). The map
\(V \mapsto v\) is a linear bijection \(\bdd{\hilb, \complex^n} \to
\complex^n \tensor \hilb\), with inverse \(V^\adjoint e_j = (e_j^\adjoint \tensor \id)v\).
\end{proof}

\begin{theorem}[Pure maps]
\label{thm:pure}
For a nonzero CP map \(\varphi\) on \(\alg = \bigoplus_k \Mn{n_k}\), the following
are equivalent:
\begin{enumerate}[label=(\roman*)]
  \item \(\varphi\) is pure;
  \item \(\pi\) is irreducible;
  \item the support projection of \(\Fphi\) is a minimal projection of \(\comm\);
  \item \(\varphi_k \ne 0\) for exactly one \(k\), and \(\operatorname{rank}C_{\varphi_k} = 1\);
  \item \(\varphi(a) = V^\adjoint a_k V\) for some \(k\) and some nonzero bounded
    \(V : \hilb \to \complex^{n_k}\).
\end{enumerate}
\end{theorem}
\begin{proof}
(i)\(\Leftrightarrow\)(ii) follows Naimark's Theorem~2. By
\cref{thm:rn}\ref{it:rn-alg}, CP maps \(\psi \le \lambda\varphi\) correspond to
positive \(T \in \pi(\alg)'\), with \(\varphi \leftrightarrow \id\). So \(\varphi\) is
pure if and only if every positive element of \(\pi(\alg)'\) is a scalar, if and
only if \(\pi(\alg)' = \complex\id\), which is irreducibility by Schur's lemma.
(ii)\(\Leftrightarrow\)(iii): \(\pi(\alg)' = P\comm P\), and this is \(\complex P\)
if and only if \(P\) is minimal in \(\comm\).
(iii)\(\Leftrightarrow\)(iv): by \cref{thm:choi} and the block decomposition of
\cref{sec:cp},
\(\comm \cong \bigoplus_k \id_{n_k} \tensor \bdd{\complex^{n_k} \tensor \hilb}\) and
\(P \cong \bigoplus_k \id_{n_k} \tensor \operatorname{supp}C_{\varphi_k}\). The minimal
projections of this algebra are \(\id_{n_k} \tensor Q\) with \(Q\) of rank one in a
single summand. (iv)\(\Leftrightarrow\)(v) is \cref{lem:kraus}.
\end{proof}

\begin{remark}
This recovers the known characterization: a CP map on \(\Mn n\) spans an extreme ray
exactly when it has Choi rank (Kraus rank) one, and in general exactly when its
minimal Stinespring representation is irreducible~\cite{Arveson1969}. In the
present approach the equivalence with the Choi rank is a statement about minimal
projections of \(\comm = \bdd\hilb \tensor \alg\op\).
\end{remark}

\subsection{Extreme points of normalized maps (Naimark \S19.4)}
\label{sec:extreme}

Naimark normalizes positive functionals by \(f(e) = 1\) and studies the extreme
points of the resulting compact convex set. For CP maps there are two natural
normalizations:
\begin{itemize}
  \item \emph{Unital (UCP)}: \(\varphi(1) = \id_{\hilb}\). By
    \cref{thm:linear}\ref{it:comm}, this is the linear condition
    \(V^\adjoint\Fphi V = \id_{\hilb}\).
  \item \emph{Trace preserving (CPTP)}: for \(\hilb = \complex^d\) and a fixed faithful
    positive functional \(\omega\) on \(\alg\), \(\trace\compose\varphi = \omega\). The
    usual choice is \(\omega = \trace_\alg \eqdef \sum_k \trace \compose p_k\).
\end{itemize}

\begin{lemma}
\label{lem:slice}
For \(G = \sum_j S_j \tensor R_{b_j} \in \comm\), the map
\(\varphi_G(a) = V^\adjoint\rho(a)GV\) of \cref{thm:linear} is
\(\varphi_G(a) = \sum_j \tau(ab_j)\,S_j\). If \(\hilb = \complex^d\), then
\(\trace\varphi_G(a) = \tau(a\,d_G)\), where
\(R_{d_G} = (\trace \tensor \id)(G)\). Hence \(\varphi\) is CPTP with respect to
\(\omega = \tau(\cdot\,d_\omega)\) if and only if
\((\trace \tensor \id)(\Fphi) = R_{d_\omega}\).
\end{lemma}
\begin{proof}
\(V^\adjoint(\id \tensor \lambda(a)R_b)V = \ip{\vect{ab}}{\Omega}\id = \tau(ab)\id\).
Take traces and use linearity of \(b \mapsto R_b\).
\end{proof}

For \(\alg = \Mn n\) and \(\omega = \trace = n\tau\), this says
\((\trace \tensor \id)\Fphi = n\id\), which is the familiar condition that the partial
trace of the Choi matrix over the output is \(\id_n\).

\begin{theorem}[Extreme points]
\label{thm:extreme}
Let \(\Lambda\) be a linear map from maps \(\alg \to \bdd\hilb\) to a space with
involution, with \(\Lambda(\psi^\sharp) = \Lambda(\psi)^\adjoint\), and let
\(\mathcal{S} = \set{\varphi \text{ CP} : \Lambda(\varphi) = \beta}\). A map
\(\varphi \in \mathcal{S}\) is an extreme point of \(\mathcal{S}\) if and only if
\[
  \pi(\alg)' \ni T \longmapsto \Lambda(\psi_T)
\]
is injective. In particular:
\begin{enumerate}[label=(\alph*)]
  \item \label{it:ext-ucp}
    (UCP, \(\Lambda(\psi) = \psi(1)\).) \(\varphi\) is extreme in
    \(\mathrm{UCP}(\alg, \bdd\hilb)\) if and only if \(T \mapsto W^\adjoint TW\) is
    injective on \(\pi(\alg)'\).
  \item \label{it:ext-tp}
    (CPTP, \(\Lambda(\psi) = \trace\compose\psi\).) \(\varphi\) is extreme among
    CPTP maps if and only if the map
    \[
      \pi(\alg)' \ni T \longmapsto (\trace \tensor \id)\br{\Fphi^{1/2}T\Fphi^{1/2}}
    \]
    is injective.
\end{enumerate}
\end{theorem}
\begin{proof}
Since \(\Lambda(\psi_{T^\adjoint}) = \Lambda(\psi_T)^\adjoint\), injectivity on
\(\pi(\alg)'\) is equivalent to injectivity on its self-adjoint part.

Suppose the map is injective and \(\varphi = \frac12(\varphi_1 + \varphi_2)\) with
\(\varphi_i \in \mathcal{S}\). Then \(\varphi_i \le 2\varphi\), so
\(\varphi_i = \psi_{T_i}\) by \cref{thm:rn}. Now
\(\Lambda(\psi_{T_1 - \id}) = \beta - \beta = 0\), so \(T_1 = \id\) and
\(\varphi_1 = \varphi\).

Conversely, suppose \(S = S^\adjoint \ne 0\) in \(\pi(\alg)'\) has
\(\Lambda(\psi_S) = 0\), scaled so that \(\norm S \le 1\). Then
\(\psi_{\id \pm S}\) are CP by \cref{thm:rn}\ref{it:rn-alg}, lie in \(\mathcal{S}\),
and average to \(\varphi\). They are distinct because \(T \mapsto \psi_T\) is
injective. Finally, (a) uses \(\psi_T(1) = W^\adjoint TW\), and (b) uses
\cref{lem:slice} with \(F_{\psi_T} = \Fphi^{1/2}T\Fphi^{1/2}\).
\end{proof}

Part~\ref{it:ext-ucp} is Arveson's extreme point
theorem~\cite[Section~1.4]{Arveson1969}. Translating through \cref{thm:choi} and
\cref{lem:kraus} gives the Choi--Kraus forms.

\begin{corollary}[Choi--Kraus characterizations]
\label{cor:extreme-kraus}
Let \(\alg = \bigoplus_k \Mn{n_k}\), \(\hilb = \complex^d\), and let
\(\varphi(a) = \sum_k\sum_{i=1}^{\kappa_k} V_{k,i}^\adjoint a_k V_{k,i}\) be a minimal
Kraus decomposition, with \(\kappa_k = \operatorname{rank}C_{\varphi_k}\).
\begin{enumerate}[label=(\alph*)]
  \item \label{it:ck-ucp}
    \(\varphi\) is extreme in \(\mathrm{UCP}(\alg, \Mn d)\) if and only if
    \(\set{V_{k,i}^\adjoint V_{k,j}}_{k,i,j}\) is linearly independent in \(\Mn d\).
    In that case \(\sum_k \kappa_k^2 \le d^2\).
  \item \label{it:ck-tp}
    \(\varphi\) is extreme among CPTP maps (\(\omega = \trace_\alg\), that is,
    \(\sum_i V_{k,i}V_{k,i}^\adjoint = \id_{n_k}\) for every \(k\)) if and only if, for
    each \(k\), \(\set{V_{k,i}V_{k,j}^\adjoint}_{i,j}\) is linearly independent in
    \(\Mn{n_k}\). In that case \(\kappa_k \le n_k\) for every \(k\).
  \item \label{it:ck-choi}
    (Choi-matrix form, \(\alg = \Mn n\), \(P = \operatorname{supp}\Choi\).) \(\varphi\)
    is extreme in UCP if and only if \(X \mapsto \sum_j (e_j^\adjoint \tensor \id)
    X (e_j \tensor \id)\) is injective on \(P(\Mn n \tensor \Mn d)P\). It is extreme
    among CPTP maps if and only if \(X \mapsto (\id \tensor \trace)(X)\) is injective
    on \(P(\Mn n \tensor \Mn d)P\).
\end{enumerate}
\end{corollary}
\begin{proof}
Since \(\dim\hilb < \infty\), \(\Fphi^{1/2}\) is invertible on \(\dil\). So
\(\set{F_{\psi_T} : T \in \pi(\alg)'} = P_{\Fphi}\comm P_{\Fphi}\), where
\(P_{\Fphi}\) is the support of \(\Fphi\). By \cref{thm:choi} and the block
decomposition, this is
\(\bigoplus_k \id_{n_k} \tensor P_k(\Mn{n_k} \tensor \Mn d)P_k\) with
\(P_k = \operatorname{supp}C_{\varphi_k} = \Span\set{v_{k,i}}_i\). By
\cref{lem:kraus}, its elements are the Choi matrices of the maps
\(\psi_s(a) = \sum_k\sum_{i,j} s^{(k)}_{ij}V_{k,i}^\adjoint a_kV_{k,j}\), with
\(s = (s^{(k)}) \in \bigoplus_k \Mn{\kappa_k}\), and \(s \mapsto \psi_s\) is
injective because the \(v_{k,i}\) are independent.

For (a), \(\psi_s(1) = \sum_{k,i,j}s^{(k)}_{ij}V_{k,i}^\adjoint V_{k,j}\) is injective in
\(s\) exactly when the family is independent, and counting dimensions gives
\(\sum_k\kappa_k^2 \le d^2\).

For (b), \(\trace\psi_s(a) = \sum_k\trace\big(a_k\sum_{i,j}s^{(k)}_{ij}V_{k,j}V_{k,i}^\adjoint\big)\)
is block diagonal in \(k\). It is injective exactly when each family is independent,
and counting gives \(\kappa_k^2 \le n_k^2\).

For (c), \(\psi_X(1)\) and \(\trace\compose\psi_X\) are the stated maps of the Choi
matrix \(X\).
\end{proof}

\begin{remark}[Comparison with the literature]
\label{rem:extreme-lit}
For \(\alg = \Mn n\), \ref{it:ck-ucp} is Choi's theorem on extreme unital CP
maps~\cite[Theorem~5]{Choi1975}. Since \(\varphi \mapsto \varphi^\adjoint\)
(Hilbert--Schmidt adjoint) exchanges CPTP maps \(\Mn n \to \Mn d\) with UCP maps
\(\Mn d \to \Mn n\), and sends Kraus operators \(V_i^\adjoint\) to \(V_i\),
\ref{it:ck-tp} is the standard criterion that a channel with Kraus operators
\(K_i\) is extreme exactly when \(\set{K_i^\adjoint K_j}\) is linearly
independent~\cite{Watrous2018}. The bounds \(\kappa \le d\) and \(\kappa \le n\) are
Choi's. The block-diagonal versions for direct sums, and the derivation from the
single operator \(\Fphi \in \bdd\hilb \tensor \alg\op\), are the parts specific to
this approach. \todo{Check theorem numbers in Choi~1975, Arveson~1969 and
Watrous~2018, and search for the direct-sum statements in the literature.}
\end{remark}

\begin{corollary}[Krein--Milman, Naimark \S19.4]
\label{cor:krein-milman}
For \(\hilb = \complex^d\), the sets of UCP and of CPTP maps on \(\alg\) are compact
and convex. So every such map is a finite convex combination of extreme points,
whose minimal Kraus decompositions satisfy the bounds of
\cref{cor:extreme-kraus}.
\end{corollary}
\begin{proof}
Both sets are closed. They are bounded because \(\norm\varphi = \norm{\varphi(1)}\)
for CP maps, and \(\norm{\varphi(1)} \le \trace\varphi(1) = \omega(1)\) in the
trace-preserving case. Apply the finite-dimensional Krein--Milman (Carath\'eodory)
theorem.
\end{proof}

\todo{Extend to the intersection UCP \(\cap\) CPTP (unital channels), where the
constraint \(\Lambda\) combines both slices. Is the criterion injectivity of
\(T \mapsto (W^\adjoint TW, (\trace \tensor \id)(\Fphi^{1/2}T\Fphi^{1/2}))\), and
does it recover the Landau--Streater conditions? Also treat extreme
\(k\)-positive maps, and quantum instruments.}

% ===========================================================================
\section{Further Directions}
\label{sec:further}

\paragraph{What is gained.}
\todo{Discuss what is gained: the dilation space is a concrete subspace of a fixed
space independent of \(\varphi\); positivity of the form is positivity of one
operator; the commutant structure explains why the Choi matrix appears; compare
with~\cite{Kasparov1980,Arveson1969,Paulsen2002}.}

\paragraph{Applications.}
\todo{Candidates to check: Kraus decompositions from a spectral decomposition of
\(\Fphi\); a Radon--Nikodym theorem for CP maps \(\psi \le \varphi\) in terms of
\(\pi(\alg)'\), which here is a compression of \(\bdd\hilb \tensor R(\alg\op)\);
quantum channels and instruments.}

\paragraph{Open questions.}
\begin{enumerate}[label=(Q\arabic*)]
  \item \emph{Finite von Neumann algebras.} For \(\alg\) a finite von Neumann
    algebra with a faithful normal tracial state, \(\Ltwo\) is a genuine
    completion, and \(\lambda(\alg)' = R(\alg\op)\) is the standard commutation
    theorem~\cite{AnantharamanPopa}. Is the \(\varphi\)-form still implemented by
    an operator affiliated with \(\bdd\hilb \bar\tensor \alg\op\)? Under what
    conditions on \(\varphi\) is it bounded?
  \item \todo{Non-tracial states (modular theory), or C\(^*\)-algebras without
    a faithful trace.}
  \item \todo{Question.}
\end{enumerate}

\bibliographystyle{plain}
\bibliography{references}

\end{document}
